\section{Rational Unification}

In this section we introduce rational terms, equation systems and the unification algorithm. The algorithm itself is inspired by the Martelli-Rossi algorithm ``UNIFY-1''~\cite{martelli1984efficient}.

\subsection{Rational Terms}

In relational (and logic) programming, we conventionally consider terms as finite expressions of the following syntax (also named Herbrand trees):
\begin{equation*}
    \begin{array}{ccl}
        \V & \text{---} & \text{variables}, \\
        \C & \text{---} & \text{constructors}, \\
        T & ::= & \mathbb{V} \mid \mathbb{C}\,(T, \dots, T). \\
    \end{array}
\end{equation*}

This notion could be generalized to consider also possibly-infinite expressions of the same syntax which we name infinite terms. While ``truly'' infinite terms like $f_1\,(f_2\,(f_3\,(\dots)))$ are not much semantically helpful in computational logic (since we cannot run programs infinitely long), it is pretty interesting to consider rational terms.

Rational term is an infinite term with the only finite set of subterms~\cite{colmerauer1982prolog} (which are subtrees). For example, $f\,(f\,(f\,(\dots)))$ is a rational term since it have only one subterm (itself).

Formally, we consider countable infinite sets $x, y, z, ... \in \V$ of variables and $f, g, h, ... \in \C$ of constructors. The latter is equipped with the function $\arity : \C \rightarrow \N$ which determines the arities of constructors and the notation $f^n \in \C$ where $f \in \C$ and $n = \arity(f)$.

We consider an endofunctor $F : \text{Set} \rightarrow \text{Set}$ and define the sets of terms $T = \mu F$ and infinite terms $T^\infty = \nu F$ as the least fixed point and the greatest fixed point~\cite{simon2006coinductive} of $F$ respectively:
\[F(X) = \V \cup \{ f^n(t_1, \dots, t_n) \mid f^n \in \C, t_i \in X \}.\]

The set of subterms of a term, $\subterms : T^\infty \rightarrow 2^{T^\infty}$, is defined in conventional way inductively:
\begin{itemize}
    \item $t \in \subterms(t)$;
    \item $t \in \subterms(t_i) \implies t \in \subterms(f^n(\dots, t_i, \dots))$.
\end{itemize}

Then, the set of rational terms is defined as $T^R = \{ t \in T^\infty \mid \subterms(t) \text{ is finite} \}$. Thus, $T \subset T^R \subset T^\infty$.

We define a substitution $\sigma \in \Sigma$ as a mapping $\V \rightarrow T^\infty$ with a finite substitution domain $\dom(\sigma) = \{ x \in \V \mid \sigma(x) \neq x \}$. Substitution application and composition operations, generality and unifier properties are defined conventionally~\cite{baader2001unification}:
\begin{itemize}
    \item Substitution application is denoted in the postfix manner: $t \sigma$;
    \item Substitution composition follows the equality: $t (\sigma_1 \diamond \sigma_2) = t\sigma_2\sigma_1$;
    \item ``More general'' relation is denoted as: $\sigma_1 \preceq \sigma_2 \iff \exists \sigma, \sigma_2 = \sigma \diamond \sigma_1$.
\end{itemize}

Also, we consider a generalization of the most general unifier (MGU) property named the minimal unifying extension. We say that $\sigma_2$ is a minimal unifying extension of $\sigma_1$ for terms $t_1, t_2 \in T^\infty$ iff:
\[ \sigma_1 \preceq \sigma_2 \land t_1 \sigma_2 = t_2 \sigma_2 \land \forall \sigma \in \Sigma. ~ \sigma_1 \preceq \sigma \land t_1 \sigma = t_2 \sigma \implies \sigma_2 \preceq \sigma. \]

Thus, $\sigma$ is MGU for terms $t_1, t_2$ iff it is a minimal unifying extension of the empty substitution.


\subsection{Equation Systems}

For convenience, we will use the notation $A^? = \{ \bot \} \cup A$ for some set $A$ and call it ``optional $A$''.

We define an equation $e \in \E$ as a pair:
\begin{itemize}
    \item a pointed finite set $X \subset \V$ named ``left-hand side'' with some variable named ``representing variable'' chosen as a base point;
    \item an optional non-variable term $t \in T^? \setminus \V$ named ``right-hand side''.
\end{itemize}

For convenience, we will use the notations $\lhs : \E \rightarrow 2^\V$, $\repr : \E \rightarrow \V$ and $\rhs : \E \rightarrow T^?$ for corresponding equation components. Also, we will use the notation $\underline{x}, y, z \equiv t$, where $\lhs(\underline{x}, y, z \equiv t) = \{ x, y, z \}$, $\repr(\underline{x}, y, z \equiv t) = x$ and $\rhs(\underline{x}, y, z \equiv t) = t$.

We define an equation system $\xi \in \Xi$ as a finite set of equations $\Xi \subset 2^\E$ with pairwise disjunctive left-hand sides. For convenience, we will use the notation $\xi, e$ for the equation system extension $\{ e \} \cup \{ e' \in \xi \mid \lhs(e) \cap \lhs(e') = \emptyset \}$.

We consider our equation definition as a representation of the special case of equation system in conventional sense. For example:
\begin{equation*}
    \begin{array}{c|c}
        \text{Conventional equation system} & \text{Our equation} \\[2pt]
        \left\{\begin{array}{l@{{}={}}l}
            x & f(z) \\
            y & x \\
            z & x \\
        \end{array}\right. & \underline{x}, y, z \equiv f(z) \\
    \end{array}
\end{equation*}

In this sense, we may see that trivial equations of the form $\underline{x} \equiv \bot$ do not equate anything so we consider them illegal, i.e., $\forall x \in \V, x \in \Xi. ~ x \equiv \bot \notin \xi$. However, we will still allow them in $\E$ for the simplicity of the following definition.

The main operation defined on equation systems is $\walk : \Xi \times \V \rightarrow \E \times T$. This operation looks up for a variable in an equation system returning a corresponding pair of equation and term. More formally:
\begin{equation*}
    \walk(\xi, x) = \begin{cases}
        (\underline{x} \equiv \bot, x) & \text{if there is no such $e \in \xi$ so $x \in \lhs(e)$}, \\
        (e, \repr(e)) & \text{if $e \in \xi, x \in \lhs(e)$ and $\rhs(e) = \bot$}, \\
        (e, \rhs(e)) & \text{otherwise}. \\
    \end{cases}
\end{equation*}

Basically, we may consider our definition of equation system as an ``enhanced'' triangular substitution~\cite{baader2001unification} w.r.t. \walk. For example:
\begin{equation*}
    \begin{array}{c|c}
        \text{Triangular substitution} & \text{Equation system} \\[2pt]
        \left\{\begin{array}{l@{{}\mapsto{}}l}
            x & f(z) \\
            y & x \\
            z & w \\
        \end{array}\right. & \left\{\begin{array}{l@{{}\equiv{}}l}
            \underline{x}, y & f(z) \\
            z, \underline{w} & \bot \\
        \end{array}\right.
    \end{array}
\end{equation*}

Indeed, we may consider equation systems as substitutions which are being applied recursively to a term. More formally, we define the equation system application as follows:
\begin{algorithmic}
    \Function{ApplyEqSys}{$\xi \in \Xi$, $t \in T^\infty$}
        \If{$x \gets t \in \V$}
            \State $(\_, t) \gets \walk(\xi, x)$
        \EndIf
        \If{$x' \gets t \in \V$}
            \State \Return $x'$
        \EndIf
        \If{$f^n(t_1, \dots, t_n) \gets t$}
            \State \Return $f^n(\Call{ApplyEqSys}{\xi, t_1}, \dots, \Call{ApplyEqSys}{\xi, t_n})$
        \EndIf
    \EndFunction
\end{algorithmic}

Under the equation system application, we may consider equation systems as a special case of substitutions. More precisely, we may prove that the application of an equation system to a rational term is always rational, too.

The second key operation of equation systems is $\union : \Xi \times \V \times \V \rightarrow \Xi \times (T \times T)^?$. Given two variables $x, y$ and an equation system $\xi$, this operation builds a new equation system from $\xi$ such that $x, y$ belong to the same equation. More formally:
\begin{algorithmic}[1]
    \Function{UnionEqSys}{$\xi \in \Xi$, $x, y \in \V$}
        \State $(e_x, \_) \gets \walk(\xi, x)$
        \State $(e_y, \_) \gets \walk(\xi, y)$
        \If{$\repr(e_x) = \repr(e_y)$}
            \State \Return $(\xi, \bot)$
            \Comment{already in the one equation or equals}
        \EndIf
        \If{$\rhs(e_x) = \bot$ \textbf{and} $\rhs(e_y) = \bot$}
            \State \Return $([\xi, \lhs(e_x) \cup \lhs(e_y) \equiv \bot], \bot)$
        \EndIf
        \If{$\rhs(e_x) = \bot$}
            \State \Return $([\xi, \lhs(e_x) \cup \lhs(e_y) \equiv \rhs(e_y)], \bot)$
        \EndIf
        \If{$\rhs(e_y) = \bot$}
            \State \Return $([\xi, \lhs(e_x) \cup \lhs(e_y) \equiv \rhs(e_x)], \bot)$
        \EndIf
        \State \Return $([\xi, \lhs(e_x) \cup \lhs(e_y) \equiv \rhs(e_x)],(\rhs(e_x), \rhs(e_y)))$
    \EndFunction
\end{algorithmic}

When we add a new equation at lines 6-12, we may choose either $\repr(e_x)$ or $\repr(e_y)$ as the representing variable. The same is true about the right-hand side of the new equation at line 12: it could be either $\rhs(e_x)$ or $\rhs(e_y)$. As we can see, the second element of \union result is intended to inform a caller about a collision in right-hand sides during the union.


\subsection{Common Part}

We define a ``common part''~\cite{martelli1984efficient} of two terms $\commonpart : T \times T \rightarrow T^?$ inductively:
\begin{equation*}
    \begin{split}
        \commonpart(x, \_) &= x \\
        \commonpart(\_, y) &= y \\
        \commonpart(f^n(l_1, \dots, l_n), f^n(r_1, \dots, r_n)) &= f^n(\commonpart(l_1, r_1), \dots, \commonpart(l_n, r_n)) \\
    \end{split}
\end{equation*}

The definition is intentionally left non-total since two terms may not have a common part (is such cases the result is $\bot$) and slightly non-deterministic (in the case of two variables) to allow optimizations. Hereafter we require the implementation of \commonpart to be a total deterministic function.


\subsection{The Unification Algorithm}

We define the unification algorithm as a recursive function $\unify : \Xi \times T \times T \rightarrow \Xi^?$:

\begin{algorithmic}[1]
    \Function{RationalUnify}{$\xi \in \Xi$, $t_1, t_2 \in T$}
        \If{$x \gets t_1 \in \V$ \textbf{and} $y \gets t_2 \in \V$}
            \State $(\xi, ts) \gets \union(\xi, x, y)$
            \If{$(t_1, t_2) \gets ts$}
                \State \Return $\Call{RationalUnify}{\xi, t_1, t_2}$
            \EndIf
            \State \Return $\xi$
        \EndIf
        \If{$x \gets t_1 \in \V$}
            \State $(e, \_) \gets \walk(\xi, x)$
            \If{$\rhs(e) = \bot$}
                \State \Return $\lhs(e) \equiv t_2, \xi$
            \EndIf
            \State $t \gets \commonpart(\rhs(e), t_2)$
            \If{$t = \bot$}
                \State \bf failure
            \EndIf
            \State $\xi \gets \lhs(e) \equiv t, \xi$
            \State \Return \Call{RationalUnify}{$\xi, \rhs(e), t_2$}
        \EndIf
        \If{$y \gets t_2 \in \V$}
            \State \Return \Call{RationalUnify}{$\xi, y, t_1$}
        \EndIf
        \If{$f^n(l_1, ..., l_n) \gets t_1$ \textbf{and} $g^m(r_1, ..., r_m) \gets t_2$}
            \If{$f^n \neq g^m$}
                \State \bf failure
            \EndIf
            \For{$i = 1, ..., n$}
                \State $\xi \gets \Call{RationalUnify}{\xi, l_i, r_i}$
            \EndFor
            \State \Return $\xi$
        \EndIf
    \EndFunction
\end{algorithmic}

We have proved the following properties of this algorithm in \Rocq\footnote{\url{https://anonymous.4open.science/r/20BC/}}:

\begin{itemize}
    \item Correctness: if the algorithm terminates with a non-$\bot$ result then the result is a minimal unifying extension of the source equation system for given terms.
    \item Completeness: if the algorithm terminates with $\bot$ result then there is no unifying extension of the source equation system for given terms.
    \item Termination: the algorithm terminates with some result on every input.
\end{itemize}
